\documentclass{article}

\IfFileExists{neurips_2018.sty}{
  \usepackage[final]{neurips_2018}
}{
  \usepackage[margin=1in]{geometry}
  \usepackage{times}
}

\usepackage[utf8]{inputenc}
\usepackage[T1]{fontenc}
\usepackage{hyperref}
\hypersetup{hidelinks}
\usepackage{url}
\usepackage{booktabs}
\usepackage{amsmath}
\usepackage{amsfonts}
\usepackage{microtype}
\usepackage{array}
\usepackage{graphicx}
\usepackage{placeins}
\usepackage{caption}

\title{Bayesian Hierarchical Logistic Regression for Telco Customer Churn}
\date{September 2026 portfolio update}

\author{%
  Justin Y. Wang\\
  Department of Applied Mathematics and Statistics\\
  Johns Hopkins University\\
  Baltimore, MD 21218\\
  \texttt{jwang518@jhu.edu}
}

\begin{document}
\maketitle

\begin{abstract}
This project studies associations with churn in the IBM Telco Customer Churn sample using a Bayesian hierarchical logistic regression with contract-specific intercepts. A hand-written Metropolis--Hastings sampler was improved by interweaving centered and non-centered updates and using a training-data covariance proposal. Four-chain rank-normalized split diagnostics support the reported training posterior. On a stratified held-out set of 1,407 customers, the Bayesian model achieved ROC-AUC 0.839, Brier score 0.138, and log loss 0.424, essentially matching ordinary logistic regression with contract fixed effects. Contract-intercept estimates were also close across models; the practical gain from partial pooling is therefore limited in this sample. The Bayesian model provides a coherent posterior description of associations, while the predictive result remains limited to one random split of a static observational data set.
\end{abstract}

\section{Introduction}

Customer churn occurs when a customer stops doing business with a company. In subscription businesses such as telecommunications, churn is especially important because acquiring new customers is often more expensive than retaining existing ones. The purpose of this project is to identify customer and contract characteristics associated with churn while using a Bayesian model that gives direct uncertainty quantification.

The main research questions are:
\begin{enumerate}
  \item Which customer-level features are associated with higher or lower churn probability?
  \item Does baseline churn risk vary meaningfully by contract type?
  \item What does a Bayesian hierarchical model add beyond a standard frequentist logistic regression baseline?
\end{enumerate}

The final question is important for the narrative of the project. The Bayesian model is not presented as a black-box classifier designed to maximize prediction accuracy. Instead, it is an interpretable probabilistic model. Its value is that it estimates uncertainty in coefficients and group-level churn risk, incorporates regularization through priors, and partially pools contract-specific intercepts.

\section{Related work}

Customer churn prediction has a large applied literature in telecommunications. Keramati et al.\ study churn prediction using several methods including logistic regression, decision trees, neural networks, support vector machines, and other data-mining approaches, emphasizing that billing, account, and service features can be predictive of churn \cite{keramati2012}. Ahmad et al.\ similarly frame churn prediction as a revenue and retention problem for telecom firms and evaluate machine-learning methods for identifying likely churners \cite{ahmad2019}. Amin et al.\ compare logistic regression and LogitBoost for telecom churn prediction, showing that logistic models remain a common interpretable baseline in this setting \cite{amin2020}. More recent churn studies often emphasize feature engineering, transformation, and model comparison across algorithms such as random forests, gradient boosting, and neural networks \cite{jeb2022}.

This project differs from much of the churn-prediction literature by focusing less on maximizing predictive accuracy and more on Bayesian inference. The model is motivated by multilevel regression and partial pooling as described by Gelman and Hill \cite{gelmanhill2007}. Contract type is treated as a grouping variable because it represents a natural segmentation of customer commitment: month-to-month, one-year, and two-year contracts. From a computational perspective, hierarchical models can mix poorly under centered parameterizations when group effects and population-level hyperparameters are strongly coupled. The non-centered parameterization used here follows standard advice for improving posterior geometry in hierarchical models \cite{betancourt2019}.

\section{Data and exploratory analysis}

The data set contains customer records for a fictional telecommunications company. The raw data have 7043 observations and 21 columns. The response variable is \texttt{Churn}, encoded as 1 for customers who churned and 0 otherwise. The only missingness issue used in this analysis is that \texttt{TotalCharges} contains 11 blank values. These rows were dropped, leaving 7032 observations. The overall churn rate after cleaning is 26.6\%.

\texttt{customerID} was removed because it is only an identifier. \texttt{TotalCharges} was converted to numeric for exploratory analysis but omitted from the regression model because it is highly collinear with tenure and monthly charges. In the cleaned data, its correlation with tenure is 0.826 and its correlation with monthly charges is 0.651. Keeping all three variables would make coefficient interpretation less stable, so the model uses tenure and monthly charges as the main continuous billing-duration variables. Both continuous predictors were standardized before fitting the Bayesian model.

\begin{table}[t]
\centering
\caption{Basic data summary after cleaning.}
\begin{tabular}{lr}
\toprule
Quantity & Value \\
\midrule
Raw observations & 7043 \\
Clean observations & 7032 \\
Dropped blank \texttt{TotalCharges} rows & 11 \\
Overall churn rate & 0.266 \\
Mean tenure & 32.42 months \\
Mean monthly charges & 64.80 \\
Corr(\texttt{tenure}, \texttt{TotalCharges}) & 0.826 \\
Corr(\texttt{MonthlyCharges}, \texttt{TotalCharges}) & 0.651 \\
\bottomrule
\end{tabular}
\end{table}

Exploratory analysis suggests several strong churn patterns. Contract type is the most visible: month-to-month customers churn at much higher rates than customers with one-year or two-year contracts. Internet service is also informative, with fiber optic customers showing higher churn than DSL or no-internet customers. Payment method is especially striking: electronic check customers have a churn rate of 45.3\%, compared with 19.2\% for mailed check, 16.7\% for automatic bank transfer, and 15.3\% for automatic credit card payment. Customers with online security or tech support also show lower observed churn rates. Figure \ref{fig:eda} summarizes these exploratory patterns.

\begin{figure}[t]
\centering
\includegraphics[width=0.95\linewidth]{figures/eda_churn_rates.pdf}
\caption{Exploratory churn rates by contract type, internet service, payment method, and support/security features. These patterns motivate the model's contract-level hierarchy and customer-level predictors.}
\label{fig:eda}
\end{figure}

\begin{table}[t]
\centering
\caption{Observed churn by payment method.}
\begin{tabular}{lrr}
\toprule
Payment method & Churn rate & $n$ \\
\midrule
Electronic check & 0.453 & 2365 \\
Mailed check & 0.192 & 1604 \\
Bank transfer (automatic) & 0.167 & 1542 \\
Credit card (automatic) & 0.153 & 1521 \\
\bottomrule
\end{tabular}
\end{table}

The EDA motivates the model structure. Contract type is used as a hierarchical grouping variable because it changes baseline churn risk. The visible differences by payment method, internet service, support features, tenure, and monthly charges directly motivate including these variables as regression predictors. Thus the model is not chosen abstractly: it is designed to formalize the strongest descriptive patterns from Figure \ref{fig:eda} while adjusting for them jointly.

\section{Model}

Let $y_i$ be the churn indicator for customer $i$, and let $g[i]$ denote the contract group for that customer. The likelihood is
\begin{align}
y_i &\sim \text{Bernoulli}(p_i),\\
\text{logit}(p_i) &= \alpha_{g[i]} + x_i^\top \beta,
\end{align}
where $x_i$ contains standardized tenure and monthly charges together with binary or one-hot encoded customer features. The contract groups are month-to-month, one-year, and two-year contracts.

The centered hierarchical prior is
\begin{align}
\alpha_g &\sim \mathcal{N}(\mu_\alpha, \tau_\alpha^2),\\
\beta_j &\sim \mathcal{N}(0,2.5^2),\\
\mu_\alpha &\sim \mathcal{N}(0,5^2),\\
\tau_\alpha &\sim \text{Half-Normal}(2.5).
\end{align}
These are weakly informative priors on the log-odds scale. The $\mathcal{N}(0,2.5^2)$ prior for fixed effects allows large but not implausibly extreme odds-ratio effects after standardizing the continuous predictors. The $\mathcal{N}(0,5^2)$ prior for $\mu_\alpha$ is broad for a logistic intercept, and the Half-Normal$(2.5)$ prior keeps the group-level standard deviation positive while mildly discouraging unrealistically large between-contract variation.
For computation, I used the equivalent non-centered parameterization
\begin{align}
z_g &\sim \mathcal{N}(0,1),\\
\alpha_g &= \mu_\alpha + \tau_\alpha z_g.
\end{align}
This parameterization leaves the statistical model unchanged. In the validation fit, non-centered updates alone mixed slowly for the hyperparameters, so centered conditional updates were interwoven with them.

\section{Computation and validation design}

After dropping 11 rows with blank \texttt{TotalCharges}, the 7,032 records were divided into a stratified training set ($n=5,625$) and a held-out test set ($n=1,407$; split seed 2026). Means and standard deviations for tenure and monthly charges were learned from the training set only. The test set was used only for final scoring and model checks. This retrospective split evaluates prediction for a current customer snapshot, not churn at a specified future horizon.

The Bayesian model used four chains of 30,000 iterations, with 10,000 warmup iterations and every second later draw retained. During warmup only, random-walk proposal scales adapted toward 0.30 acceptance. The fixed-effect proposal used a covariance estimate from a frequentist logistic fit to the training data. To improve mixing for the contract hierarchy, each sweep combined non-centered Metropolis updates with a centered conditional draw for $\mu_\alpha$ and a centered Metropolis update for $\log\tau_\alpha$, holding the contract intercepts fixed. These updates target the same posterior. Final block acceptance rates ranged from 0.284 to 0.324.

Diagnostics were recalculated from separate chains with rank-normalized split $\hat R$ and approximate chain-aware bulk and tail effective sample sizes. Every parameter in the primary fit had $\hat R<1.009$; the smallest approximate bulk ESS was 366 for the month-to-month intercept. Figure \ref{fig:trace} shows two representative traces. This is a substantial improvement over the original course notebook, whose concatenated-chain ESS and unsplit $\hat R$ should not be used as convergence evidence.

\begin{figure}[!htbp]
\centering
\includegraphics[width=0.75\linewidth]{figures/mcmc_trace_validation.png}
\caption{Post-warmup traces from the four validation chains for a fixed effect and the contract hyper-mean. Complete diagnostics are saved in \texttt{results/mcmc\_diagnostics.csv}.}
\label{fig:trace}
\end{figure}

\section{Posterior results}

Table \ref{tab:posterior} summarizes training-set posteriors. Longer tenure, online security, tech support, and no internet service were associated with lower churn odds; higher monthly charges, paperless billing, fiber service, and electronic-check payment were associated with higher odds. The interval for dependents included zero in this training fit, so its direction is not established. These are adjusted associations, not effects of changing a customer's services.

\begin{table}[!htbp]
\centering
\caption{Selected training-set Bayesian fixed effects. Continuous predictors are standardized.}
\label{tab:posterior}
\begin{tabular}{lrr}
\toprule
Predictor & Mean log odds & 95\% CrI \\
\midrule
Tenure & -0.798 & $[-0.925,-0.677]$ \\
Monthly charges & 0.310 & $[0.096,0.529]$ \\
Senior citizen & 0.277 & $[0.090,0.459]$ \\
Dependents & -0.145 & $[-0.339,0.053]$ \\
Paperless billing & 0.400 & $[0.234,0.571]$ \\
Online security & -0.378 & $[-0.577,-0.187]$ \\
Tech support & -0.436 & $[-0.623,-0.248]$ \\
Fiber optic internet & 0.508 & $[0.206,0.805]$ \\
No internet service & -0.763 & $[-1.102,-0.427]$ \\
Electronic check & 0.386 & $[0.180,0.592]$ \\
\bottomrule
\end{tabular}
\end{table}

\begin{figure}[!htbp]
\centering
\includegraphics[width=0.81\linewidth]{figures/posterior_fixed_effects_validation.png}
\caption{Training posterior means and 95\% credible intervals for selected fixed effects.}
\label{fig:fixed}
\end{figure}

The contract intercepts retain the expected order, but these intercepts describe a particular reference profile, not each group's observed churn rate. The frequentist fixed-effect intercepts were within 0.02 log-odds units of the corresponding Bayesian means. With three well-populated groups, this analysis does not demonstrate a material estimation advantage from partial pooling.

\begin{center}
\captionof{table}{Training contract intercepts and frequentist comparison.}
\centering
\begin{tabular}{lrrr}
\toprule
Contract & Bayesian mean & 95\% CrI & Frequentist \\
\midrule
Month-to-month & -1.338 & $[-1.630,-1.048]$ & -1.323 \\
One year & -1.952 & $[-2.276,-1.627]$ & -1.940 \\
Two year & -2.649 & $[-3.089,-2.210]$ & -2.666 \\
\bottomrule
\end{tabular}
\end{center}

A sensitivity run changed the Half-Normal prior scale on $\tau_\alpha$ from 2.5 to 1.0. The posterior mean of $\tau_\alpha$ changed from 1.262 to 0.871, while each contract-intercept mean moved by less than 0.01. Thus the group effects are stable under this check, but the population spread remains prior-sensitive with only three groups. The alternative run's largest rank-split $\hat R$ was 1.014, so the sensitivity estimates are approximate.

\section{Held-out prediction and model checks}

The frequentist baseline used the same customer predictors and contract fixed effects. Table \ref{tab:heldout} reports test-set results alongside a constant-probability baseline set to training prevalence. Both logistic models improved substantially over that baseline and were effectively tied with each other. A one-customer difference in accuracy at threshold 0.5 does not establish a meaningful winner.

\begin{table}[!htbp]
\centering
\caption{Held-out prediction on 1,407 customers. AP is average precision.}
\label{tab:heldout}
\begin{tabular}{lrrrr}
\toprule
Model & Accuracy & ROC-AUC & AP & Brier / log loss \\
\midrule
Prevalence baseline & 0.7342 & 0.5000 & 0.2658 & 0.1952 / 0.5791 \\
Frequentist logistic & 0.8038 & 0.8391 & 0.6537 & 0.1380 / 0.4236 \\
Bayesian hierarchical & 0.8045 & 0.8390 & 0.6537 & 0.1380 / 0.4235 \\
\bottomrule
\end{tabular}
\end{table}

Figure \ref{fig:calibration} displays held-out calibration by prediction decile. To check the probability model more directly, binary churn outcomes were simulated from 1,200 posterior draws at the test-set covariates. Observed churn rates for the full test set and each contract group fell within the corresponding 95\% replicated-outcome intervals. This is a limited check: some individual calibration deciles differ visibly from the diagonal.

\begin{figure}[!htbp]
\centering
\includegraphics[width=0.57\linewidth]{figures/heldout_calibration.pdf}
\caption{Held-out calibration by prediction decile for the Bayesian model.}
\label{fig:calibration}
\end{figure}

\section{Discussion}

The updated analysis supports several associations with churn and shows that both logistic models discriminate similarly on the held-out customers. The hierarchical prior offers a principled posterior for contract intercepts, but its practical shrinkage is small here. Logistic regression is an appropriately strong baseline, and frequentist confidence intervals also quantify uncertainty; the Bayesian contribution is the joint posterior and directly stated posterior probabilities under explicit priors.

The data are observational, so these coefficients do not estimate the effect of an intervention. Tenure and service fields also require a current-customer prediction point; a signup model would need different features and a future churn horizon. This is a retrospective split after the original full-sample course analysis informed the model specification, so it is less independent than a test set reserved before EDA. It does not test future-period drift or transport to another company. The 0.5 threshold is illustrative, not an optimized retention policy. Finally, uncertainty in the hyperparameter $\tau_\alpha$ depends on its prior because there are only three contract groups.

\section{Conclusion}

A custom Bayesian hierarchical logistic model can be fit with credible chain diagnostics and evaluated honestly against a matched logistic baseline. The held-out results are effectively tied: both achieve ROC-AUC about 0.839 and log loss about 0.424. The model's value is an interpretable posterior analysis of associations, while predictive superiority and a practical benefit from partial pooling are not established by this data set.

\FloatBarrier

\begin{thebibliography}{9}

\bibitem[Ahmad et al.(2019)]{ahmad2019}
Ahmad, A. K., Jafar, A., and Aljoumaa, K. (2019).
\newblock Customer churn prediction in telecom using machine learning in big data platform.
\newblock \emph{Journal of Big Data}, 6(28).

\bibitem[Amin et al.(2020)]{amin2020}
Amin, A., Anwar, S., Adnan, A., Nawaz, M., Howard, N., Qadir, J., Hawalah, A., and Hussain, A. (2020).
\newblock Churn prediction in telecommunication using logistic regression and LogitBoost.
\newblock \emph{Procedia Computer Science}, 167, 101--112.

\bibitem[Betancourt(2019)]{betancourt2019}
Betancourt, M. (2019).
\newblock Hierarchical modeling.
\newblock \url{https://betanalpha.github.io/assets/case_studies/hierarchical_modeling.html}.

\bibitem[Gelman and Hill(2007)]{gelmanhill2007}
Gelman, A. and Hill, J. (2007).
\newblock \emph{Data Analysis Using Regression and Multilevel/Hierarchical Models}.
\newblock Cambridge University Press.

\bibitem[Jeb et al.(2022)]{jeb2022}
Jeb, S. N., Lasheng, Y., Aluvalu, R., and Dasig, D. (2022).
\newblock A novel customer churn prediction model for the telecommunication industry using data transformation methods and feature selection.
\newblock \emph{PLOS ONE}, 17(12):e0278095.

\bibitem[Keramati et al.(2012)]{keramati2012}
Keramati, A., Jafari-Marandi, R., Aliannejadi, M., Ahmadian, I., Mozaffari, M., and Abbasi, U. (2012).
\newblock Customer churn prediction in telecommunications.
\newblock \emph{Expert Systems with Applications}, 39(1), 1414--1425.

\end{thebibliography}

\end{document}
